% Generated draft: requires mathematical and source review.
\documentclass[11pt]{article}
\usepackage[margin=25mm]{geometry}
\usepackage{fontspec}
\setmainfont{Latin Modern Roman}
\usepackage{amsmath,amssymb,mathtools,unicode-math}
\setmathfont{Latin Modern Math}
\usepackage{enumitem,needspace,xurl,hyperref}
\setlength{\parindent}{0pt}
\setlength{\parskip}{5pt}
\setlength{\emergencystretch}{3em}
\newcommand{\N}{\mathbb{N}}
\newcommand{\Z}{\mathbb{Z}}
\newcommand{\Q}{\mathbb{Q}}
\newcommand{\R}{\mathbb{R}}
\newcommand{\C}{\mathbb{C}}
\newcommand{\F}{\mathbb{F}}
\newcommand{\A}{\mathbb{A}}
\newcommand{\PP}{\mathbb{P}}
\newcommand{\E}{\mathbb{E}}
\newcommand{\eps}{\varepsilon}
\newcommand{\dd}{\,\mathrm{d}}
\begin{document}
\section*{UnsolvedMath ID: 700982}
\textbf{Source problem code:} catalogue entry
\subsection{Definitions and mathematical statement}
Let \(G\) be a group. For \(g\in G\), let \(\operatorname{ord}(g)\) be the least positive integer \(m\) such that \(g^m=1_G\), and put \(\operatorname{ord}(g)=\infty\) if no such integer exists. Define
\[
\pi_e(G)=\{\operatorname{ord}(g):g\in G\}.
\]
For a positive integer \(n\), call \(G\) an \(OC_n\)-group if
\[
\pi_e(G)=\{1,2,\ldots,n\}.
\]
Thus every element has finite order at most \(n\), and every integer from \(1\) through \(n\) occurs as an element order. A group is locally finite if every subgroup generated by finitely many elements is finite.

The first question, proposed by W. J. Shi, is whether
\[
\forall n\in\Z_{\geq1}\ \forall G,\qquad
\pi_e(G)=\{1,2,\ldots,n\}
\ \Longrightarrow\ G\text{ is locally finite}.
\]
The second question asks whether infinite \(OC_n\)-groups exist for \(n\geq7\). The supplied wording leaves the quantification over \(n\) unresolved. Its existential reading is
\[
\exists n\in\Z_{\geq7}\ \exists G,\qquad
G\text{ is infinite}\quad\text{and}\quad
\pi_e(G)=\{1,2,\ldots,n\}.
\]
Its reading as an existence question for every such parameter is
\[
\forall n\in\Z_{\geq7}\ \exists G,\qquad
G\text{ is infinite}\quad\text{and}\quad
\pi_e(G)=\{1,2,\ldots,n\}.
\]
These are distinct formulations; the supplied record does not explicitly select between them.
\subsection{Short English statement}
An \(OC_n\)-group has exactly the element orders \(1,2,\ldots,n\). The problem asks whether every such group is locally finite, and whether infinite examples exist for \(n\geq7\). The second question does not explicitly distinguish existence for some \(n\geq7\) from existence for every such \(n\).
\subsection{Sources}
\begin{itemize}
\item Local UnsolvedMath registry, permanent ID \(700982\), source code \texttt{catalogue entry}, title \emph{Kourovka Notebook Problem 13.64}. Supplied upstream revision: \texttt{d371ee917ac9647fbb06cd6fc529deec6f3dda82}. The registry labels the problem open; this is supplied source metadata, not an independently verified assessment.
\item \emph{The Kourovka Notebook (Unsolved Problems in Group Theory)}, original problem 13.64, 13th issue (1995); proposed by W. J. Shi. The supplied record also identifies the 21st issue edition (2026), arXiv:1401.0300v48. Its collection description names E. I. Khukhro and V. D. Mazurov as editors. These bibliographic details are reproduced from the supplied record without independent source verification.
\end{itemize}
\end{document}
